% explainer-source-hash: 861400b897d7b1fa48553bd9cae0133347dd1552fda7a27409d9a04edad8c953
% explainer-source-commit: aaa9d3ff8bf9a423477b01fafdb5b66d2a478f1c
% explainer-generated: 2026-07-16
% Regenerate with /explain-all (see WIRING.md). Do not edit this stamp by hand:
% it is how the toolchain knows whether this document still matches the code.
\documentclass[11pt,a4paper]{article}
\input{preamble}

% ---- local helper: one small 3x3 board with a win line highlighted -------
% #1 = x shift, #2 = comma list of cell indices to highlight, #3 = caption
\newcommand{\miniboard}[3]{%
  \begin{scope}[xshift=#1]
    \foreach \i in {#2} {
      \pgfmathtruncatemacro{\rr}{int(\i/3)}
      \pgfmathtruncatemacro{\cc}{mod(\i,3)}
      \fill[accent!35] (\cc, 2-\rr) rectangle ++(1,1);
    }
    \draw[linegrey, thin] (0,0) grid (3,3);
    \draw[inkblue, thin] (0,0) rectangle (3,3);
    \node[font=\small\ttfamily, text=inkblue] at (1.5,-0.55) {#3};
  \end{scope}
}

\begin{document}

\begin{titlepage}
\centering
\vspace*{3cm}
{\Huge\bfseries\color{inkblue} Tic Tac Toe (tkinter GUI)\par}
\vspace{0.8cm}
{\Large\color{accent} Deep Dive\par}
\vspace{1.5cm}
{\large A clickable 3x3 game window whose rules and unbeatable\\
minimax opponent live in a module that has never heard of a window.\par}
\vspace{2cm}
{\normalsize Two Python files, 495 lines of source, zero third-party dependencies.\par}
\vfill
{\small Written from the source code. Every claim below was read out of the\\
repository or produced by running it.\par}
\vspace{1cm}
\end{titlepage}

\tableofcontents
\newpage

% =========================================================================
\section{What this project is}

This is a Tic Tac Toe game you play in a real window with a mouse. You click a
square, your mark appears in it, and the program decides whether anyone has won.
It has two modes: two people taking turns at the same keyboard, or one person
against a computer opponent that cannot be beaten.

That last claim is the interesting part of the project, and it is the reason the
codebase is shaped the way it is. ``Unbeatable'' is not decoration here. It is a
specific, mechanically checkable property, and the repository contains the check.

\begin{keyidea}{The one idea to take away}
The game's \emph{rules} and the game's \emph{window} are kept in separate files
that know nothing about each other. \file{game.py} can decide who won without a
screen existing; \file{main.py} draws pixels and never decides anything about the
rules. This is why the rules and the AI can be tested by a script with no display
attached, which in turn is why the ``unbeatable'' claim could be verified at all.
\end{keyidea}

\subsection{The whole repository}

\begin{figure}[H]
\centering
\begin{forest} dirtree
[tic-tac-toe-gui
  [game.py \qquad {\rmfamily\itshape 154 lines: rules + AI, no I/O}]
  [main.py \qquad {\rmfamily\itshape 341 lines: the tkinter window}]
  [tests/
    [test\_ai.py \qquad {\rmfamily\itshape 6 unit tests on AI decisions}]
    [verify\_unbeatable.py \qquad {\rmfamily\itshape standalone proof script}]
  ]
  [docs/
    [screenshots/app.png]
    [explainers/ \qquad {\rmfamily\itshape this document}]
  ]
  [README.md]
  [requirements.txt \qquad {\rmfamily\itshape deliberately empty of packages}]
  [LICENSE \qquad {\rmfamily\itshape PolyForm Strict 1.0.0}]
  [.gitignore]
]
\end{forest}
\caption{Every tracked file. There are only two source files.}
\end{figure}

\begin{jargon}{Module}
In Python, a module is simply one \file{.py} file. Writing \code{import game} in
another file runs \file{game.py} once and hands you back everything it defined
(its functions and variables) so you can call them. ``Keeping \file{game.py} free
of tkinter'' means: nothing inside that file mentions windows, so importing it
never tries to open one.
\end{jargon}

\begin{jargon}{tkinter}
The graphical toolkit that ships inside Python itself. It gives you objects for
windows, buttons and text labels, and an event loop (see below) that waits for
the user to do something. It needs no installation on Windows or macOS; on Debian
and Ubuntu it is a separate system package, \code{python3-tk}. It is the only
reason this project needs anything beyond bare Python, and it is not a
third-party dependency: it is part of the standard library.
\end{jargon}

\begin{jargon}{Standard library}
The set of modules that arrive with Python itself and need no installation. This
project uses exactly three: \code{tkinter} (the window), \code{math} (for the
values \code{-inf} and \code{inf}), and \code{random} plus \code{unittest} in the
tests. \file{requirements.txt} therefore lists no packages at all; it contains
only a comment saying so.
\end{jargon}

\subsection{Why it exists, per the README}

The README states the origin plainly: this began as a ``Day 83'' console practice
script that printed a numbered board and read \code{input()} from the terminal,
and that script had a real bug. Its win check compared cells only against the
string \code{" X "}, so player O could complete a line and the game would carry on
as though nothing had happened. The rewrite keeps the original exercise (represent
a 3x3 board and detect a win across all eight lines without a library), fixes the
bug at its root, and replaces the terminal loop with an actual window.

That history is worth holding onto, because the fix is visible in the current
code's shape. Section~\ref{sec:winlines} shows how.

% =========================================================================
\section{Architecture: two layers and a one-way arrow}

\begin{figure}[H]
\centering
\resizebox{\textwidth}{!}{
\begin{tikzpicture}[node distance=8mm]

  \node[box, minimum width=34mm] (user) {The person\\(mouse clicks)};
  \node[box, right=16mm of user, minimum width=44mm, minimum height=30mm] (gui)
    {\textbf{main.py}\\\code{TicTacToeApp}\\[2pt]
     \scriptsize widgets, colors, score,\\\scriptsize turn order, the 400\,ms delay};
  \node[box, right=16mm of gui, minimum width=44mm, minimum height=30mm] (logic)
    {\textbf{game.py}\\[2pt]
     \scriptsize \code{make\_move}, \code{check\_winner},\\
     \scriptsize \code{is\_draw}, \code{minimax},\\
     \scriptsize \code{best\_ai\_move}};
  \node[store, right=16mm of logic, minimum height=22mm] (board) {\code{board}\\9-cell\\list};

  \node[extbox, above=10mm of gui, minimum width=44mm] (tk) {tkinter\\\scriptsize (standard library)};
  \node[extbox, below=12mm of logic, minimum width=44mm] (tests)
    {\textbf{tests/}\\\code{\scriptsize test\_ai.py}\\\code{\scriptsize verify\_unbeatable.py}};

  \draw[flow] (user) -- node[lbl, above] {click} (gui);
  \draw[flow] (gui) -- node[lbl, above] {calls\\functions} (logic);
  \draw[flow] (logic) -- node[lbl, above] {reads\\and writes} (board);
  \draw[dflow] (logic) to[bend right=18] node[lbl, below] {returns winner,\\line, or move} (gui);
  \draw[dflow] (gui) to[bend right=18] node[lbl, below] {redraws\\cell} (user);
  \draw[flow] (tk) -- (gui);
  \draw[flow] (tests) -- node[lbl, right] {import\\directly} (logic);

  \begin{scope}[on background layer]
    \node[fit=(gui)(tk), draw=linegrey, dashed, rounded corners, inner sep=5mm,
          label={[font=\small\itshape, text=black!60]below:presentation layer}] {};
    \node[fit=(logic)(board)(tests), draw=linegrey, dashed, rounded corners, inner sep=5mm,
          label={[font=\small\itshape, text=black!60]above:pure logic layer}] {};
  \end{scope}
\end{tikzpicture}}
\caption{The dependency arrow points one way only. \file{main.py} imports
\file{game.py}; \file{game.py} imports nothing but \code{math}. The tests reach
past the window entirely and talk to the logic directly, which is only possible
because that arrow never points back.}
\label{fig:arch}
\end{figure}

\begin{keyidea}{Why the one-way arrow matters}
If \file{game.py} imported tkinter, then \file{tests/verify\_unbeatable.py} could
not run: importing the rules would try to touch a display, and on a machine with
no screen (a build server, a container) it would fail outright. The 642 exhaustive
games in Section~\ref{sec:verify} were playable precisely because no window was
ever involved. The separation is not tidiness for its own sake; it is what makes
the project's headline claim checkable.
\end{keyidea}

% =========================================================================
\section{\file{game.py}: the rules}

\subsection{How a board is represented}

The board is a flat Python list of nine strings, not a 3x3 nested list. Each entry
is \code{"X"}, \code{"O"}, or \code{""} (the constant \code{EMPTY}). The indices
run left to right, top to bottom:

\begin{figure}[H]
\centering
\begin{tikzpicture}[scale=1.0]
  \foreach \i in {0,...,8} {
    \pgfmathtruncatemacro{\rr}{int(\i/3)}
    \pgfmathtruncatemacro{\cc}{mod(\i,3)}
    \draw[linegrey] (\cc, 2-\rr) rectangle ++(1,1);
    \node[font=\large\ttfamily, text=inkblue] at (\cc+0.5, 2.5-\rr) {\i};
  }
  \draw[inkblue] (0,0) rectangle (3,3);
  \node[font=\small\itshape, text=black!60, align=left, anchor=west] at (3.6,1.5)
    {\code{board[4]} is the centre.\\\code{board = ["X","","O", ...]}\\is one list of nine strings.};
\end{tikzpicture}
\caption{The flat indexing scheme, copied from the module's own docstring.}
\end{figure}

\begin{jargon}{Flat list vs nested list}
A nested list would be \code{[["X","",""],["","O",""],["","",""]]}, addressed as
\code{board[row][col]}. A flat list is \code{["X","","","","O","","","",""]},
addressed as \code{board[4]}. Flat wins here because every win line becomes three
plain integers, the win table becomes eight tuples of numbers, and the minimax
search can place and unplace a move with a single assignment. Nothing in this
project ever wants a whole row as an object, so the extra structure would only
cost.
\end{jargon}

\subsection{\code{WIN\_LINES} and the bug that shaped it}
\label{sec:winlines}

There are exactly eight ways to win. The module writes them down once, as data:

\begin{lstlisting}[language=Python, caption={game.py, the win table}]
WIN_LINES = [
    (0, 1, 2), (3, 4, 5), (6, 7, 8),  # rows
    (0, 3, 6), (1, 4, 7), (2, 5, 8),  # columns
    (0, 4, 8), (2, 4, 6),             # diagonals
]
\end{lstlisting}

\begin{figure}[H]
\centering
\resizebox{\textwidth}{!}{
\begin{tikzpicture}
  \miniboard{0cm}{0,1,2}{(0,1,2)}
  \miniboard{3.6cm}{3,4,5}{(3,4,5)}
  \miniboard{7.2cm}{6,7,8}{(6,7,8)}
  \miniboard{10.8cm}{0,3,6}{(0,3,6)}
  \miniboard{14.4cm}{1,4,7}{(1,4,7)}
  \miniboard{18.0cm}{2,5,8}{(2,5,8)}
  \miniboard{21.6cm}{0,4,8}{(0,4,8)}
  \miniboard{25.2cm}{2,4,6}{(2,4,6)}

  \draw[decorate, decoration={brace, amplitude=5pt}, inkblue]
    (0,3.35) -- (9.6,3.35) node[midway, above=4pt, font=\small, text=inkblue] {3 rows};
  \draw[decorate, decoration={brace, amplitude=5pt}, inkblue]
    (10.8,3.35) -- (20.4,3.35) node[midway, above=4pt, font=\small, text=inkblue] {3 columns};
  \draw[decorate, decoration={brace, amplitude=5pt}, inkblue]
    (21.6,3.35) -- (27.6,3.35) node[midway, above=4pt, font=\small, text=inkblue] {2 diagonals};
\end{tikzpicture}}
\caption{All eight entries of \code{WIN\_LINES}, one board each. Every one is
three indices into the flat list, nothing more.}
\end{figure}

The original console version, per the README, hardcoded eight \code{or}-chained
comparisons and every one of them tested for \code{" X "}. Writing the same eight
comparisons again for O would have doubled a block that was already hard to read,
and it is exactly the kind of block in which one missing case hides. The rewrite
sidesteps that by turning the eight lines into a list and looping:

\begin{lstlisting}[language=Python, caption={game.py, check\_winner}]
def check_winner(board):
    for line in WIN_LINES:
        a, b, c = line
        if board[a] != EMPTY and board[a] == board[b] == board[c]:
            return board[a], line
    return None
\end{lstlisting}

The mark is never named in the comparison. The test is ``these three cells hold
the same thing, and that thing is not empty'', which is true for X and O with one
code path. The bug becomes unrepresentable rather than merely fixed.

\begin{keyidea}{Data instead of branches}
The fix is not ``add the missing O check''. It is ``stop writing the check per
mark''. Eight near-identical \code{if} statements cannot be audited by eye; eight
tuples in a list can be checked against the picture above in seconds. Turning
repeated code into a table of data is the move here, and it is the project's
central lesson according to its own README.
\end{keyidea}

The return value is worth noting: \code{check\_winner} gives back the winning mark
\emph{and the line}, as a tuple, or \code{None}. The line is not needed to decide
the game; it is needed so the GUI can paint those three cells gold. The rules layer
returns it because the presentation layer will want it, but it returns it as three
plain integers, not as a colour or a widget. That is the boundary being respected.

\subsection{The rest of the rules}

\begin{table}[H]
\centering
\small
\begin{tabular}{@{}p{3.7cm}p{4.2cm}p{6.6cm}@{}}
\toprule
\textbf{Function} & \textbf{Returns} & \textbf{Behaviour} \\
\midrule
\code{new\_board()} & a fresh 9-item list & \code{[EMPTY] * 9}. \\
\code{make\_move(board, index, mark)} & \code{True} / \code{False} &
  Writes \code{mark} into the cell if the index is in 0..8 \emph{and} the cell is
  empty; otherwise changes nothing and returns \code{False}. The caller decides
  what a rejection means. \\
\code{check\_winner(board)} & \code{(mark, line)} or \code{None} &
  Scans the eight lines, in table order, and stops at the first complete one. \\
\code{is\_draw(board)} & \code{True} / \code{False} &
  Board full \emph{and} no winner. Both halves are required: a full board with a
  winner is a win, not a draw. \\
\code{is\_game\_over(board)} & \code{True} / \code{False} &
  Winner exists, or draw. \\
\code{available\_moves(board)} & list of indices & Every empty cell, in board order. \\
\code{other\_mark(mark)} & \code{"X"} or \code{"O"} &
  \code{"O" if mark == "X" else "X"}. \\
\bottomrule
\end{tabular}
\caption{The seven public helpers. None of them raise; all of them are total
functions over any 9-item list.}
\end{table}

\code{make\_move} is the one worth pausing on, because of the order of its test:

\begin{lstlisting}[language=Python, caption={game.py, make\_move}]
if index < 0 or index > 8 or board[index] != EMPTY:
    return False
\end{lstlisting}

The range check comes first for a reason: Python's \code{or} short-circuits, so if
\code{index} is 12, the expression stops before it ever evaluates
\code{board[12]}, which would raise \code{IndexError}. Reordering these three
conditions would turn a clean \code{False} into a crash. The guard is correct as
written and is fragile to rearrangement, which is the kind of thing worth knowing
before touching it.

\begin{honest}{\code{make\_move} silently accepts a nonsense mark}
Nothing validates that \code{mark} is \code{"X"} or \code{"O"}.
\code{make\_move(board, 0, "banana")} returns \code{True} and puts the string
\code{"banana"} on the board, and \code{check\_winner} would happily declare
``banana'' the winner of a completed line. In practice no caller can do this: both
\file{main.py} and the tests only ever pass marks that came from
\code{other\_mark} or a hardcoded literal. So this is a latent sharp edge, not a
live bug. It is listed here because a reader auditing \file{game.py} in isolation
should know its guarantees end at the index check.
\end{honest}

\begin{honest}{\code{other\_mark} treats everything that is not X as O's opposite}
\code{other\_mark("banana")} returns \code{"X"}, because the implementation is a
single conditional on \code{mark == "X"} with no third branch. Same category as
above: unreachable given the actual callers, but not a function you would want to
expose to untrusted input without a guard.
\end{honest}

% =========================================================================
\section{The AI: minimax}

This is the part of the project that justifies its existence. Everything above is
bookkeeping; this is an actual algorithm.

\subsection{The problem, stated precisely}

Given a board, and knowing it is my turn, which cell should I play so that I never
lose? Not ``probably do well'' and not ``use good heuristics like take the centre'':
never lose, against any opponent, ever.

Tic Tac Toe is small enough that this question can be answered by brute force. From
an empty board there are at most $9! = 362{,}880$ orderings of the nine cells, and
far fewer real games because most end before the board fills. A computer can look
at all of them faster than a human can blink.

\begin{jargon}{Game tree}
A picture of every way a game can unfold. The root is the position right now.
Each branch is one legal move. Each leaf is a finished game (someone won, or the
board filled). To ``search the game tree'' is to walk every branch to its leaf and
see what happened.
\end{jargon}

\begin{jargon}{Recursion}
A function that calls itself on a smaller version of the same problem.
\code{minimax} asks ``how good is this board for me?'' by placing each legal move
and asking the same question about each resulting board, until it hits a board
where the answer is obvious because the game is over. That obvious case is the
\emph{base case}, and without one the function would call itself forever.
\end{jargon}

\subsection{Scoring a finished game}

At a leaf, the answer is not a guess. It is a fact:

\begin{lstlisting}[language=Python, caption={game.py, \_score\_terminal}]
def _score_terminal(board, ai_mark, depth):
    result = check_winner(board)
    if result is None:
        return 0
    winner = result[0]
    return (10 - depth) if winner == ai_mark else (depth - 10)
\end{lstlisting}

A win for the AI is positive, a loss is negative, a draw is zero. Scores are always
from the AI's point of view, no matter whose turn produced the leaf. That single
convention is what lets one number flow all the way back up the tree without ever
needing to be flipped.

\begin{jargon}{depth}
How many moves deep in the search this board is. The board handed to
\code{best\_ai\_move} is depth 0; each move the search tries adds one.
\end{jargon}

The \code{depth} term is a tie-breaker. A win at depth 1 scores $10 - 1 = 9$; the
same win at depth 5 scores $10 - 5 = 5$. Both are wins, but the shallower one
scores higher, so given two winning paths the AI takes the shorter one. It will
close out a game it has already won instead of shuffling around the board first.
Symmetrically, a loss at depth 7 scores $7 - 10 = -3$, which is less bad than a
loss at depth 1 at $-9$, so a cornered AI stalls as long as possible.

\begin{honest}{The losing half of the depth term is dead code in practice}
The module's own docstring concedes this: the AI prefers ``losing later (if a loss
were ever forced, which minimax guarantees it never is against best play)''. On a
3x3 board the second player can always force at least a draw, so the AI is never in
a genuinely lost position when it is asked to move from a legal game state. The
$depth - 10$ branch \emph{does} get evaluated constantly during the search (the
search must consider losing lines in order to avoid them), but the tie-break it
implements between two losses can never influence the root decision in a real game.
It is correct, harmless, and by the code's own admission unreachable in effect.
\end{honest}

\begin{honest}{The magic number 10 is never explained, and does not need to be}
It only needs to exceed the maximum depth (9) so that scores never cross zero and
change sign: with 10, the worst win still scores $10 - 9 = 1 > 0$ and the mildest
loss scores $9 - 10 = -1 < 0$. Had the constant been 8, a win at depth 9 would
score $-1$ and the AI would read some wins as losses. Nothing in the code, the
comments or the README states this constraint. It is a load-bearing invariant held
together by a bare literal, and there is no test that would catch it being lowered.
(The relationship between 10 and the depth bound is my reading of the arithmetic,
labelled \emph{inferred}; the code does not say it.)
\end{honest}

\subsection{Minimax proper}

The AI wants a high score; the opponent, playing well, wants a low one. So the
search alternates: on the AI's turns take the maximum of the children, on the
opponent's turns take the minimum.

\begin{figure}[H]
\centering
\resizebox{0.94\textwidth}{!}{
\begin{tikzpicture}[level distance=20mm, sibling distance=26mm,
  every node/.style={box, minimum width=11mm, minimum height=7mm, font=\small}]

  \node (root) {\textbf{9}}
    child { node (a) {\textbf{9}}
      child { node {9} }
      child { node {5} }
    }
    child { node (b) {0}
      child { node {0} }
      child { node {7} }
    }
    child { node (c) {-8}
      child { node {-8} }
      child { node {3} }
    };

  \node[lbl, left=16mm of root, text=inkblue, font=\small\bfseries] (l1) {MAX\\(AI to move)};
  \node[lbl, left=6mm of a, text=inkblue, font=\small\bfseries, xshift=-16mm] (l2) {MIN\\(opponent)};
  \node[lbl, below left=0mm and 6mm of a, text=inkblue, font=\small\bfseries, xshift=-16mm, yshift=-14mm] (l3) {leaves\\(game over)};

  \node[lbl, right=4mm of root, text=warnred] {takes the biggest child: 9};
  \node[lbl, above right=1mm and 2mm of b, text=warnred] {takes the smallest: 0};
\end{tikzpicture}}
\caption{Minimax on a toy tree. The leaf scores come from
\code{\_score\_terminal}. Each MIN node assumes the opponent will pick the worst
outcome for the AI, so it reports the minimum of its children. The root then picks
the branch reporting the highest number. Crucially the AI does not hope the
opponent blunders: it assumes perfect defence and still finds a line it is happy
with. That pessimism is why the guarantee holds against \emph{any} opponent, not
just a good one.}
\label{fig:minimax}
\end{figure}

\begin{pseudocode}{minimax(board, current\_mark, ai\_mark, depth, alpha, beta)}
\Step{if the game on `board` is over:}
\Indent{return score\_terminal(board, ai\_mark, depth), None}
\Step{}
\Step{moves := every empty cell}
\Step{best\_move := moves[0]}
\Step{}
\Step{if current\_mark is the AI:}
\Indent{best\_score := -infinity}
\Indent{for each move in moves:}
\Indent{\hspace*{2em}place current\_mark on board[move]}
\Indent{\hspace*{2em}score := minimax(board, other(current\_mark), ai\_mark, depth+1, alpha, beta)}
\Indent{\hspace*{2em}undo: board[move] := EMPTY}
\Indent{\hspace*{2em}if score > best\_score: best\_score, best\_move := score, move}
\Indent{\hspace*{2em}alpha := max(alpha, best\_score)}
\Indent{\hspace*{2em}if beta <= alpha: break \quad // prune}
\Step{else:  // the opponent}
\Indent{... identical, but minimising, tracking beta, same prune test}
\Step{}
\Step{return best\_score, best\_move}
\end{pseudocode}

Three details in the real code deserve naming.

\textbf{The board is mutated and restored, not copied.} The loop writes
\code{board[move] = current\_mark}, recurses, then writes
\code{board[move] = EMPTY} again. A copy per node would be simpler to reason about
but would allocate a new list hundreds of thousands of times. The undo makes the
function's contract exact: the docstring promises ``\code{board} is left unchanged
on return'', and it is, because every write is paired with its own undo on every
path, including the pruning \code{break} (the undo happens before the comparison
that can break).

\begin{honest}{The undo pattern is correct but unguarded}
There is no \code{try}/\code{finally} around the place-recurse-undo sequence. If
anything inside the recursion raised, the board would be left with a phantom mark
on it and the caller's board would be silently corrupted. Nothing in this call
chain can raise (it is arithmetic and list indexing on a validated structure), so
this is theoretical. But the correctness of the ``unchanged on return'' promise
rests on ``no exceptions'', not on the language enforcing it.
\end{honest}

\textbf{\code{best\_move} is seeded with \code{moves[0]}.} If every move scored
identically, the function still returns a legal move rather than \code{None}. This
matters because the strict comparison \code{if score > best\_score} never fires on a
tie, so without the seed a board where all children tie would return no move at all.

\textbf{The AI is deterministic.} There is no randomness anywhere in
\code{minimax}. Ties are broken by enumeration order, which is board order, which
means the first-listed cell achieving the best score wins. On an empty board this
produces index 0, the top-left corner, every single time. I confirmed this by
running it: \code{best\_ai\_move(new\_board(), "X")} returns \code{0}, and the full
empty-board search completes in about 0.055 seconds on this machine.

\begin{honest}{A comment in the tests claims the opening is the centre; it is not}
\file{tests/test\_ai.py} contains the comment ``a center-first opening is the
classic optimal choice and is what this specific minimax tie-breaking $\dots$
should produce''. It does not. The code returns corner 0, because 0 is enumerated
before 4 and both score 0. The comment then partially catches itself in the next
sentence and the assertion is loosened to
\code{assertIn(move, [0, 2, 4, 6, 8])}, which passes. So the test is green and its
prose is wrong. Verified by running the function directly. Both a corner and the
centre are game-theoretically fine openings, so nothing is broken; the comment is
just stale reasoning left in the file.
\end{honest}

\subsection{Alpha-beta pruning, and what it is not}

\begin{jargon}{Alpha-beta pruning}
An optimisation that skips branches which cannot change the answer. \code{alpha}
is the best score the maximiser has already guaranteed itself somewhere else;
\code{beta} is the best the minimiser has guaranteed. When \code{beta <= alpha},
whatever remains in this branch is unreachable under sensible play by both sides,
so the loop stops early.
\end{jargon}

The intuition, stated as a scenario: you are the maximiser and you have already
found a line worth 5. You start exploring another line, and its first reply already
drops you to 3. The opponent chooses that reply, so this whole line is worth at
most 3. It cannot beat the 5 you already have. There is no reason to look at the
rest of it, no matter how good it might get, because you will never be allowed to
go there.

\begin{keyidea}{Pruning is not why the AI is unbeatable}
This distinction is easy to blur and the README explicitly calls it out as a place
the author had to correct himself. The unbeatability comes \emph{entirely} from
exhaustive minimax search. Delete both \code{alpha} and \code{beta} from the code
and the AI is exactly as unbeatable, just slower. Conversely, bolt alpha-beta onto
a broken search and you get a fast broken search. Pruning changes how much of an
already-decided tree gets visited; it never changes the decision. The two claims
are independent and should never be merged into ``the AI is unbeatable because of
alpha-beta minimax''.
\end{keyidea}

\begin{honest}{Pruning is present but unmeasured, and unnecessary at this size}
The code has no benchmark comparing pruned to unpruned, and there is no test that
would notice if the pruning were wrong (a bug in the cutoff would show up as a
worse move, which the existing tests would catch only by luck). Its practical value
here is close to zero: the full empty-board search takes 0.055 seconds \emph{with}
pruning, and 3x3 Tic Tac Toe is small enough that the unpruned version would also
finish faster than a human notices. Alpha-beta is in this codebase because it is
the correct companion to minimax and worth being able to explain, not because the
project needed the speed. Saying otherwise would be inventing a rationale the code
does not support.
\end{honest}

\subsection{\code{best\_ai\_move}: the entire public surface of the AI}

\begin{lstlisting}[language=Python, caption={game.py, best\_ai\_move}]
def best_ai_move(board, ai_mark):
    if is_game_over(board):
        return None
    _, move = minimax(board, current_mark=ai_mark, ai_mark=ai_mark)
    return move
\end{lstlisting}

Four lines. \code{current\_mark} and \code{ai\_mark} are the same at the root, because
the root is by definition the AI's turn; they diverge one level down and stay
alternating from there. The score is thrown away, because the caller does not want
to know how confident the AI is, only where to click.

This is the whole interface \file{main.py} uses. Give it a board and a mark, get an
index back. The GUI has no idea a tree search happened.

% =========================================================================
\section{\file{main.py}: the window}

\subsection{The widget tree}

\begin{jargon}{Widget}
One on-screen thing: a window, a button, a text label, a radio button. In tkinter
you create widgets, tell each one who its parent is, and then \emph{pack} it,
which asks the parent to find it a place. Widgets form a tree, and the whole
window is that tree drawn.
\end{jargon}

\begin{figure}[H]
\centering
\begin{forest} dirtree
[{root \quad {\rmfamily\itshape the Tk window; title "Tic Tac Toe"; not resizable}}
  [{header \quad {\rmfamily\itshape Frame}}
    [{status\_label \quad {\rmfamily\itshape "Player X's turn" / "AI is thinking..." / the result}}]
  ]
  [{mode\_frame \quad {\rmfamily\itshape Frame}}
    [{Radiobutton "Two player" \quad {\rmfamily\itshape writes self.mode}}]
    [{Radiobutton "Vs AI" \quad {\rmfamily\itshape writes self.mode}}]
  ]
  [{mark\_frame \quad {\rmfamily\itshape Frame; hidden unless mode is vs\_ai}}
    [{Label "You play:"}]
    [{Radiobutton "X (first)" \quad {\rmfamily\itshape writes self.human\_mark}}]
    [{Radiobutton "O (second)"}]
  ]
  [{grid\_frame \quad {\rmfamily\itshape Frame}}
    [{buttons 0..8 \quad {\rmfamily\itshape nine Buttons; .grid(row=i//3, col=i\%3)}}]
  ]
  [{controls \quad {\rmfamily\itshape Frame}}
    [{play\_again\_btn \quad {\rmfamily\itshape calls reset\_board}}]
  ]
  [{score\_frame \quad {\rmfamily\itshape Frame}}
    [{score\_label \quad {\rmfamily\itshape "X wins: 0 \ \ O wins: 0 \ \ Draws: 0"}}]
  ]
]
\end{forest}
\caption{Built once by \code{\_build\_layout}. Every frame is packed vertically;
only the nine cells use \code{.grid}.}
\end{figure}

\begin{jargon}{pack and grid}
tkinter's two layout managers. \code{pack} stacks children in order (top to bottom
by default); \code{grid} places them in rows and columns. They cannot be mixed
\emph{inside the same parent}, but different parents may use different ones, which
is exactly what happens here: the frames are packed into the window, and the nine
buttons are gridded into \code{grid\_frame}.
\end{jargon}

The nine buttons are created in a loop, and the line that wires them up is the one
Python beginners get wrong:

\begin{lstlisting}[language=Python, caption={main.py, the cell loop}]
for i in range(9):
    row, col = divmod(i, 3)
    btn = tk.Button(
        grid_frame, ...,
        command=lambda idx=i: self.on_cell_click(idx),
    )
    btn.grid(row=row, column=col, padx=4, pady=4, ipadx=8, ipady=8)
    self.buttons.append(btn)
\end{lstlisting}

\begin{keyidea}{Why \code{lambda idx=i} and not \code{lambda: ...(i)}}
Python closures capture the \emph{variable}, not its value at the time the lambda
was written. Writing \code{lambda: self.on\_cell\_click(i)} would give all nine
buttons a function reading the same \code{i}, which by the time anyone clicks has
finished the loop and equals 8. Every cell would play the bottom-right square. The
\code{idx=i} default argument is evaluated once, when the lambda is created, and
freezes the value. This is a genuine bug avoided, not a style choice.
\end{keyidea}

\code{divmod(i, 3)} returns the quotient and remainder together, so index 5 becomes
row 1, column 2. It is the same flat-to-2D mapping the board uses, which is why the
button list and the board list can share a single index with no translation
anywhere in the file.

\subsection{The state the GUI owns}

\begin{table}[H]
\centering
\small
\begin{tabular}{@{}p{4.4cm}p{3.2cm}p{7.0cm}@{}}
\toprule
\textbf{Attribute} & \textbf{Type} & \textbf{Meaning} \\
\midrule
\code{self.board} & 9-item list & The only game state. Owned here, operated on by \file{game.py}. \\
\code{self.current\_player} & \code{"X"} / \code{"O"} & Whose turn. Always starts at \code{"X"}. \\
\code{self.game\_over} & bool & Set once a round ends; blocks all clicks. \\
\code{self.ai\_thinking} & bool & True during the 400\,ms pause; blocks all clicks. \\
\code{self.mode} & \code{tk.StringVar} & \code{"two\_player"} or \code{"vs\_ai"}. \\
\code{self.human\_mark} & \code{tk.StringVar} & \code{"X"} or \code{"O"}; only read in AI mode. \\
\code{self.scores\_two\_player} & dict & Keys \code{"X"}, \code{"O"}, \code{"draws"}. \\
\code{self.scores\_vs\_ai} & dict & Keys \code{"human"}, \code{"ai"}, \code{"draws"}. \\
\bottomrule
\end{tabular}
\caption{Note what is absent: no copy of the win lines, no ``did X win'' flag, no
duplicate rule logic. The GUI stores what the GUI needs and asks \file{game.py}
everything else.}
\end{table}

\begin{jargon}{tk.StringVar}
A tkinter box holding a string that widgets can watch. Both radio buttons in a
group are handed the same \code{StringVar} and different \code{value}s; clicking
one writes its value into the box and visually deselects the other. Reading it is
\code{self.mode.get()}, never \code{self.mode}.
\end{jargon}

The AI's mark is not stored at all. It is a computed property:

\begin{lstlisting}[language=Python, caption={main.py, ai\_mark}]
@property
def ai_mark(self):
    return other_mark(self.human_mark.get())
\end{lstlisting}

A stored copy could drift out of sync with the radio button. A property cannot,
because there is nothing to keep in sync: it is derived on every read. Small
decision, but it is the same instinct as \code{WIN\_LINES}, which is to hold one
source of truth and compute the rest.

\subsection{Two independent score dictionaries}

Two-player scores are keyed by mark (\code{"X"}, \code{"O"}); AI-mode scores are
keyed by role (\code{"human"}, \code{"ai"}). This is not an inconsistency, it is the
right call: in AI mode the player can switch between X and O between rounds, so a
mark-keyed tally would scatter one person's wins across two counters. Keying by
role survives the switch. In two-player mode there is no role, only a mark, so the
mark is the natural key. Two different questions, two different keys.

\begin{honest}{The scores vanish when the window closes}
They are plain in-memory dictionaries, never written anywhere. Closing the app
resets everything to zero. The README's ``What I'd do differently'' names this
explicitly as a known gap, so it is a stated limit rather than an oversight.
\end{honest}

\subsection{The click path, end to end}

\begin{figure}[H]
\centering
\resizebox{\textwidth}{!}{
\begin{tikzpicture}[node distance=6mm and 12mm]

  \node[box] (click) {Human clicks\\cell $i$};
  \node[dec, right=of click] (guard) {game over\\or AI\\thinking?};
  \node[box, below=9mm of guard, fill=white, draw=linegrey] (drop1) {ignore};
  \node[dec, right=of guard] (turn) {vs AI and\\not human's\\turn?};
  \node[dec, right=of turn] (mv) {make\_move\\accepted?};
  \node[box, right=of mv] (render) {\_render\_cell($i$)\\\scriptsize draw mark, disable button};
  \node[dec, below=13mm of render] (over) {\_finish\_turn\_if\_over};
  \node[box, left=of over] (end) {record score,\\highlight line,\\set game\_over};
  \node[box, below=11mm of over] (switch) {\_switch\_turn()};
  \node[dec, left=of switch] (isai) {vs AI and\\AI's turn?};
  \node[box, left=of isai] (sched) {\_schedule\_ai\_move\\\scriptsize ai\_thinking = True\\\scriptsize root.after(400, ...)};
  \node[box, below=9mm of isai] (status) {\_update\_status()};
  \node[box, below=10mm of sched] (aimove) {\_make\_ai\_move\\\scriptsize best\_ai\_move -> make\_move\\\scriptsize -> render -> finish or switch};

  \draw[flow] (click) -- (guard);
  \draw[flow] (guard) -- node[lbl, above] {no} (turn);
  \draw[flow] (guard) -- node[lbl, right] {yes} (drop1);
  \draw[flow] (turn) -- node[lbl, above] {no} (mv);
  \draw[flow] (turn) |- node[lbl, near start, right] {yes} (drop1);
  \draw[flow] (mv) -- node[lbl, above] {yes} (render);
  \draw[flow] (mv) |- node[lbl, near start, right] {no} (drop1);
  \draw[flow] (render) -- (over);
  \draw[flow] (over) -- node[lbl, above] {game ended} (end);
  \draw[flow] (over) -- node[lbl, right] {still going} (switch);
  \draw[flow] (switch) -- (isai);
  \draw[flow] (isai) -- node[lbl, above] {yes} (sched);
  \draw[flow] (isai) -- node[lbl, right] {no} (status);
  \draw[dflow] (sched) -- node[lbl, right] {after 400\,ms} (aimove);
\end{tikzpicture}}
\caption{\code{on\_cell\_click}, the only entry point a human has into the game.
Three guards reject the click before anything happens; the rest is a straight line
through \file{game.py} and back out to the screen.}
\label{fig:click}
\end{figure}

The three guards deserve enumerating, because between them they cover every way a
click can be illegitimate:

\begin{enumerate}
\item \code{self.game\_over or self.ai\_thinking}: the round is finished, or the AI's
  400\,ms timer is pending. Both are ``the board is not yours right now''.
\item \code{mode == "vs\_ai" and current\_player != human\_mark.get()}: it is the
  AI's turn. The comment in the source calls these ``stray human clicks''.
\item \code{not make\_move(...)}: the cell is already taken. Note the GUI does not
  check this itself; it lets \file{game.py} decide and reads the boolean. The rules
  layer is the authority even on the trivial question.
\end{enumerate}

Guard 3 is doubly covered: filled cells have their button set to
\code{state="disabled"} by \code{\_render\_cell}, so tkinter will not even fire the
command. The \code{make\_move} check is a second line of defence that would still
hold if the disabling were ever removed.

\subsection{The 400\,ms delay and the event loop}

\begin{jargon}{Event loop}
\code{root.mainloop()} is a loop that never returns until the window closes. It
waits for events (clicks, key presses, timers) and calls your functions in
response. Everything your code does happens \emph{inside} a callback from that
loop, and while your callback runs, the loop is blocked and the window is frozen.
\end{jargon}

\begin{jargon}{root.after(ms, fn)}
Asks the event loop to call \code{fn} once, about \code{ms} milliseconds from now,
and returns immediately. It is not a thread and not a sleep: the window stays
responsive in the meantime because control goes straight back to the loop.
\end{jargon}

The AI's move is computed in about 0.055 seconds in the worst case (empty board,
verified by running it), so it would appear instantly. The code inserts a
deliberate pause:

\begin{lstlisting}[language=Python, caption={main.py, scheduling the AI}]
AI_MOVE_DELAY_MS = 400

def _schedule_ai_move(self):
    self.ai_thinking = True
    self._update_status(thinking=True)
    self.root.after(AI_MOVE_DELAY_MS, self._make_ai_move)
\end{lstlisting}

The source comment is explicit that this is cosmetic: ``The move itself is computed
instantly; this delay is purely a UX touch so the AI doesn't feel like it's
snapping the board shut.'' Worth being precise about what it is not: it is not
waiting for a computation, and it is not \code{time.sleep(0.4)}, which would freeze
the entire window for 400\,ms because the event loop would be blocked inside the
callback. Using \code{after} keeps the window alive; the \code{ai\_thinking} flag is
what stops the human from using that liveness to click during the pause.

\begin{honest}{The \code{after} handle is discarded, and ``Play again'' can make the AI move twice}
\code{root.after} returns an id that \code{root.after\_cancel(id)} could use. The
return value is thrown away, and \code{after\_cancel} appears nowhere in the file
(verified by grep: \code{after} occurs once, as a call).

The consequence is reachable through the UI. Suppose in AI mode the human moves,
the AI's timer starts, and 100\,ms later the human clicks ``Play again''.
\code{reset\_board} clears the board and sets \code{ai\_thinking = False} and
\code{game\_over = False}, but it cannot stop the pending timer. If the human then
plays again, a \emph{second} timer is scheduled while the first is still live. Both
fire. Each one calls \code{\_make\_ai\_move}, which re-checks only
\code{self.game\_over} (false) and not whose turn it actually is, so the AI plays a
mark, the turn switches to the human, and then the stale callback plays another AI
mark out of turn.

The fix is one line at the schedule site (keep the id) and one in
\code{reset\_board} (cancel it), or a turn check inside \code{\_make\_ai\_move}. I
could not drive this through the real GUI on this machine (no display server and no
\code{Xvfb} available), so the failure is traced from the source rather than
observed: \textbf{labelled inferred}. The discarded return value and the absent
\code{after\_cancel} are verified facts; the resulting double move is a reading of
the control flow.
\end{honest}

\subsection{The round as a state machine}

\begin{figure}[H]
\centering
\resizebox{0.96\textwidth}{!}{
\begin{tikzpicture}[node distance=16mm and 26mm]
  \node[box] (human) {\textbf{HUMAN\_TURN}\\\scriptsize clicks accepted};
  \node[box, right=of human] (think) {\textbf{AI\_THINKING}\\\scriptsize ai\_thinking = True\\\scriptsize clicks ignored};
  \node[box, below=of think] (aiturn) {\textbf{AI\_MOVED}\\\scriptsize mark placed, cell disabled};
  \node[box, below=of human] (done) {\textbf{ROUND\_OVER}\\\scriptsize game\_over = True\\\scriptsize score recorded, line gold};

  \draw[flow] (human) -- node[lbl, above] {click, game continues,\\vs AI mode} (think);
  \draw[dflow] (think) -- node[lbl, right] {after 400\,ms:\\\_make\_ai\_move} (aiturn);
  \draw[flow] (aiturn) -- node[lbl, above] {game continues} (human);
  \draw[flow] (human) -- node[lbl, left] {click completes\\a line or fills\\the board} (done);
  \draw[flow] (aiturn) to[bend left=22] node[lbl, below] {AI's move ends it} (done);
  \draw[flow] (done) to[out=160, in=200, looseness=3.5] node[lbl, left] {Play again:\\reset\_board} (human);
  \draw[flow] (human) to[out=110, in=70, looseness=4] node[lbl, above] {click, game continues,\\two-player mode\\(turn switches, same state)} (human);
\end{tikzpicture}}
\caption{The lifecycle of one round. In two-player mode the AI\_THINKING and
AI\_MOVED states are never entered; the self-loop on HUMAN\_TURN is the whole game.
\code{reset\_board} is the only edge back out of ROUND\_OVER, and in AI mode where
the AI moves first it fires \code{\_schedule\_ai\_move} on the way, landing directly
in AI\_THINKING rather than HUMAN\_TURN.}
\end{figure}

\subsection{\code{\_finish\_turn\_if\_over}, and the argument nobody reads}

\begin{lstlisting}[language=Python, caption={main.py, end-of-round handling}]
def _finish_turn_if_over(self, mark_just_played):
    result = check_winner(self.board)
    if result is not None:
        winner, line = result
        self.game_over = True
        self._record_result(winner=winner)
        self._highlight_line(line)
        self._update_status(winner=winner)
        self._update_score_label()
        return True

    if is_draw(self.board):
        self.game_over = True
        self._record_result(draw=True)
        self._update_status(draw=True)
        self._update_score_label()
        return True

    return False
\end{lstlisting}

One function handles both endings, for both players, in both modes, and both call
sites (the human's click and the AI's move) use it identically. That is the payoff
of the layer split: the GUI's end-of-round handling is five lines of bookkeeping
because the hard question was answered elsewhere.

\begin{honest}{\code{mark\_just\_played} is dead}
The parameter is never read in the body. The function asks \code{check\_winner}
who won rather than trusting the caller's word for it, which is arguably the better
design, but then the argument should have been deleted. Both call sites dutifully
pass one (\code{just\_played} from the click handler, \code{self.ai\_mark} from the
AI path), and both values are discarded. Verified by reading the body: the only
names it touches are \code{self.board} and its own locals.
\end{honest}

\subsection{Keeping the marks coloured after a cell is disabled}

\begin{lstlisting}[language=Python, caption={main.py, \_render\_cell}]
self.buttons[index].config(
    text=mark, fg=color, state="disabled", disabledforeground=color
)
\end{lstlisting}

tkinter greys out a disabled button's text by default, so the blue X and the red O
would both fade to the same grey the instant they were placed, and the board would
lose its most basic visual signal. Setting \code{disabledforeground} to the same
colour as \code{fg} defeats that. The README lists this as one of the project's
real challenges, and it is a good example of the category: not algorithmically
hard, just a toolkit default that silently does the wrong thing.

\code{\_highlight\_line} does the same trick in reverse, repainting the three
winning cells gold with dark text, again setting \code{disabledforeground} because
those buttons are already disabled by the time anyone has won.

\subsection{Showing and hiding the mark chooser}

\begin{lstlisting}[language=Python, caption={main.py, mode switching}]
def _on_mode_change(self):
    if self.mode.get() == "vs_ai":
        self.mark_frame.pack(padx=16, pady=(0, 8), fill="x",
                             before=self._grid_frame())
    else:
        self.mark_frame.pack_forget()
    self._update_score_label()
    self.reset_board()

def _grid_frame(self):
    return self.buttons[0].master
\end{lstlisting}

\code{pack\_forget} removes a widget from the layout without destroying it, so it
can be packed again later with its state intact. But re-packing appends to the
bottom of the parent by default, which would drop the mark chooser underneath the
board. The \code{before=} argument pins it back above the grid.
\code{\_grid\_frame()} recovers that grid frame by asking button 0 who its parent
is, since the frame was a local variable in \code{\_build\_layout} and was never
stored on \code{self}.

\begin{honest}{\code{\_grid\_frame} is a workaround for a one-line omission}
Storing \code{self.grid\_frame = grid\_frame} in \code{\_build\_layout} would have
made the helper, and its explanatory comment, unnecessary. Reaching through
\code{buttons[0].master} works and is correct, but it re-derives a fact the code
already had and then needs a comment to justify itself. Minor, and cosmetic.
\end{honest}

Note the side effect: \code{\_on\_mode\_change} is wired to \emph{all four} radio
buttons, including the two mark selectors. So changing your mark mid-game resets
the board, with no confirmation. Scores survive (they are not touched), but the
position is gone. Whether that is a bug or the intent is not stated anywhere; it
is at least defensible, since changing sides mid-game has no meaning.

% =========================================================================
\section{Tests and the unbeatable claim}
\label{sec:verify}

\subsection{\file{tests/test\_ai.py}: six unit tests}

\begin{jargon}{Unit test}
A small function that sets up one specific situation, runs one piece of code, and
asserts the answer is what it should be. Run them all and you learn whether any of
those situations broke. Python's \code{unittest} module, in the standard library,
provides the machinery.
\end{jargon}

Both test files begin with the same three lines:

\begin{lstlisting}[language=Python]
import os, sys
sys.path.insert(0, os.path.dirname(os.path.dirname(os.path.abspath(__file__))))
from game import best_ai_move, minimax, new_board
\end{lstlisting}

This walks up from the test file to the project root and puts it on Python's import
path, so \code{import game} resolves whether you run
\code{python3 tests/test\_ai.py} directly or
\code{python3 -m unittest discover -s tests} from the root. There is no
\file{\_\_init\_\_.py} and no packaging; this is the small-project way, and it works.

\begin{table}[H]
\centering
\small
\begin{tabular}{@{}p{5.9cm}p{8.7cm}@{}}
\toprule
\textbf{Test} & \textbf{What it pins down} \\
\midrule
\code{takes\_immediate\_winning\_move} & Board \code{X X \_ / O O \_ / \_ \_ \_}, AI is X: must play 2 and win now. \\
\code{blocks\_immediate\_human\_win} & Board \code{O O \_ / X \_ \_ / \_ \_ \_}, AI is X: must play 2 to block. \\
\code{prefers\_win\_over\_block\_}\allowbreak\code{when\_both\_available} & Claims to test win-over-block. See the finding below. \\
\code{empty\_board\_optimal\_first\_move} & Empty board scores 0 (a draw with perfect play) and the move is a corner or the centre. \\
\code{returns\_none\_on\_full\_board} & \code{best\_ai\_move} on a full board returns \code{None}. \\
\code{takes\_center\_when\_it\_is\_}\allowbreak\code{the\_only\_strong\_reply} & X in a corner, AI is O: must answer 4, the centre. \\
\bottomrule
\end{tabular}
\caption{All six pass. Verified: \code{Ran 6 tests ... OK} in 0.047 seconds.}
\end{table}

\begin{honest}{Test 3 is a copy of test 1 and tests nothing new}
\code{test\_prefers\_win\_over\_block\_when\_both\_available} builds the board
\code{["X","X","","O","O","","","",""]} and asserts the move is 2. That is
byte-for-byte the same board and the same assertion as
\code{test\_takes\_immediate\_winning\_move}. The comment above it is a visible
train of thought that never landed: it starts to describe one layout, notices the
win cell and the block cell coincide, says ``use a layout where winning and
blocking are genuinely different cells'', starts a third idea, abandons it with
``so build a cleaner case instead:'' and then builds the original board again.

So the suite has six tests but five distinct cases, and the specific behaviour this
test names (choosing your own win over blocking theirs when they are different
cells) is not covered by anything. The comment even flags the trap and the code
walks into it anyway. Verified by comparing the two literals directly.
\end{honest}

\subsection{\file{tests/verify\_unbeatable.py}: the actual proof}

This is not a unit test and does not pretend to be. It is a standalone script that
attacks the ``unbeatable'' claim three different ways, each for both starting
orders, and exits non-zero with the losing board if the AI ever loses once.

\begin{figure}[H]
\centering
\resizebox{0.95\textwidth}{!}{
\begin{tikzpicture}[node distance=9mm and 14mm]
  \node[box, minimum width=42mm] (ex) {\textbf{1. Exhaustive}\\\scriptsize human branches into\\\scriptsize \emph{every} legal move;\\\scriptsize AI replies with minimax};
  \node[box, minimum width=42mm, right=of ex] (rand) {\textbf{2. Random}\\\scriptsize 1000 games per order,\\\scriptsize human plays uniformly\\\scriptsize at random (seed 42)};
  \node[box, minimum width=42mm, right=of rand] (opt) {\textbf{3. Optimal}\\\scriptsize human also runs minimax;\\\scriptsize theory says both orders\\\scriptsize must draw};
  \node[box, below=14mm of rand, minimum width=52mm] (v) {\textbf{verdict}\\\scriptsize any AI loss anywhere\\\scriptsize $\Rightarrow$ print the board, exit 1};

  \node[lbl, below=1mm of ex, text width=42mm, align=center]
    {569 games (human first)\\73 games (AI first)\\\textbf{0 losses}};
  \node[lbl, below=1mm of opt, text width=42mm, align=center]
    {draw, draw\\\textbf{0 losses}};

  \draw[flow] (ex) -- (v);
  \draw[flow] (rand) -- (v);
  \draw[flow] (opt) -- (v);
\end{tikzpicture}}
\caption{The three checks. Numbers shown are from an actual run on this machine,
reproduced below.}
\end{figure}

\begin{keyidea}{Why check 1 is the one that counts}
Checks 2 and 3 are supporting evidence. Random play (check 2) can only ever fail to
find a counterexample; it samples, so a loss that needs one exact 5-move sequence
could hide from a million random games. Check 3 tests a single line of play.

Check 1 is different in kind. It does not sample: at every human turn it recurses
into \emph{every} legal move, so it visits every distinct game any human strategy
whatsoever could produce against this AI. If the AI could lose, a leaf of that
traversal is that loss, and the script finds it. That is the difference between
``we did not observe a loss'' and ``a loss does not exist''. The README's ``What I
learned'' names this distinction explicitly, and it is the sharpest thinking in the
project.
\end{keyidea}

The recursion in check 1 uses the same place-recurse-undo pattern as
\code{minimax}, on its own board, and asymmetrically: the AI's turn produces one
child (minimax is deterministic, so there is exactly one AI reply per position),
while the human's turn produces one child per empty cell. That asymmetry is why the
counts are so small. 569 and 73 are not the number of Tic Tac Toe games; they are
the number of games reachable \emph{against this particular deterministic
opponent}, and the AI-first figure is smaller because the AI's fixed opening
removes a whole layer of human choice before the human ever moves.

\subsubsection*{Verified output}

I ran the script. Every figure the README states is reproduced exactly:

\begin{lstlisting}[caption={python3 tests/verify\_unbeatable.py, run on this machine, 57 seconds}]
=== 1. Exhaustive search over every human move sequence ===
  [human moves first] 569 distinct games explored -> AI won: 386, draws: 183, AI lost: 0
  [AI moves first] 73 distinct games explored -> AI won: 71, draws: 2, AI lost: 0

=== 2. 1000 games, human plays uniformly random legal moves ===
  [human moves first] 1000 games -> AI won: 789, draws: 211, AI lost: 0
  [AI moves first] 1000 games -> AI won: 996, draws: 4, AI lost: 0

=== 3. Both sides play optimally (human also uses minimax) ===
  [human moves first] outcome: draw
  [AI moves first] outcome: draw

RESULT: the AI never lost a single game across all checks above.
\end{lstlisting}

Check 2 is seeded (\code{random.Random(42)}), so those counts are reproducible
rather than merely typical, which is why the README can quote 789 and 211 as fixed
numbers.

\begin{honest}{What the verification does not cover}
It exercises \file{game.py} only. Not one line of \file{main.py} is tested by
anything, anywhere in this repository. The AI is proven; the window is not. Every
GUI behaviour in this document (the guards, the delay, the highlight, the score
dictionaries, the stale-timer bug above) rests on reading the source, because there
is no display on this machine and no \code{Xvfb} to fake one. The project's testing
stops precisely at the layer boundary, which is honest about where the interesting
logic is, but it does mean the double-move bug in \code{\_schedule\_ai\_move} could
never have been caught by the suite as it stands.
\end{honest}

\begin{honest}{The screenshot in \file{docs/} predates the AI}
\file{docs/screenshots/app.png} shows the status label, the 3x3 grid, ``Play again''
and an \code{X wins / O wins / Draws} score line, but there are no ``Two player''
and ``Vs AI'' radio buttons, and no mark chooser. \code{\_build\_layout} packs
\code{mode\_frame} with both radios unconditionally, so the current program cannot
produce that window. The image documents an earlier build. It is not referenced
from the README, so nothing states a false claim; the file is simply stale.
\end{honest}

% =========================================================================
\section{Dependencies, licence, and what is deliberately absent}

\file{requirements.txt} contains no packages at all, only a comment explaining that
there are none and noting the \code{python3-tk} caveat on Linux and the Python 3.6+
floor. The floor is real: the code uses f-strings (\code{f"Player \{winner\} wins!"}),
which arrived in 3.6.

The three imports across the whole project are \code{math} (only for
\code{math.inf}, the initial alpha and beta), \code{tkinter}, and \code{random} plus
\code{unittest} in the tests. Every one ships with Python.

\begin{jargon}{math.inf}
Positive infinity as a float. It seeds the maximiser's \code{best\_score} at
\code{-inf} so that the very first real score, however bad, is an improvement. This
avoids needing a separate ``have I looked at anything yet'' flag.
\end{jargon}

The licence is PolyForm Strict 1.0.0, copyright 2026 Salman Adnan. Strict means
exactly what it says: the code may be read but not used, modified or redistributed.
This is a portfolio artefact, not an open-source offering, and the licence matches
that intent rather than contradicting it.

\file{.gitignore} covers the usual Python debris (\code{\_\_pycache\_\_/},
\code{*.pyc}, \code{.pytest\_cache/}, \code{.venv/}, \code{venv/}) and
\code{.DS\_Store}. Note \code{.pytest\_cache/} is listed although the project uses
\code{unittest}, not pytest; harmless, and evidently a template line.

% =========================================================================
\section{Honest summary}

What this project does well is a single thing done consistently: it draws a line
between rules and pixels and never crosses it. That line is not a stylistic
preference here. It is load-bearing, because it is what let a script play 642
exhaustive games and 4000 random ones with no window in sight, and therefore what
turned ``unbeatable'' from a boast into a number. The README's own conclusion, that
a claim is only worth writing if something behind it would catch the claim being
false, is demonstrated by the repository rather than asserted by it.

The AI is a correct, textbook minimax with alpha-beta, and the project is careful
in exactly the place most write-ups are careless: it refuses to credit pruning for
the guarantee, which belongs to the exhaustive search alone.

Against that, the flaws are real and worth stating without softening:

\begin{itemize}
\item \textbf{A reachable GUI bug.} The discarded \code{after} handle lets ``Play
  again'' during the AI's pause leave a live timer that can make the AI move twice
  (inferred from source; not observed, no display available).
\item \textbf{A duplicate test.} Six tests, five cases. The one behaviour named in
  test 3's own title is not covered by anything, and its comment shows the author
  spotting the problem and then not fixing it.
\item \textbf{A wrong comment.} The test suite claims the AI opens in the centre. It
  opens at corner 0. Verified by running it.
\item \textbf{A stale screenshot.} \file{docs/screenshots/app.png} shows a build
  with no AI mode.
\item \textbf{Zero coverage of \file{main.py}.} Every GUI claim in this document is
  read, not run.
\item \textbf{Small latent sharp edges.} \code{make\_move} accepts any string as a
  mark; \code{other\_mark} has no third branch; the constant 10 in
  \code{\_score\_terminal} silently depends on exceeding the maximum depth.
\end{itemize}

None of these touch the headline claim. The AI is unbeatable, it is verified
exhaustively, and I reproduced the verification. The defects live in the layer
nobody tested, which is exactly where one would predict them.

\begin{table}[H]
\centering
\small
\begin{tabular}{@{}p{8.4cm}p{6.2cm}@{}}
\toprule
\textbf{Claim in this document} & \textbf{Basis} \\
\midrule
All README verification numbers (569, 73, 789, 211, 996, 4, zero losses) & \textbf{Verified}: script run, output reproduced above \\
Six unit tests pass in 0.047\,s & \textbf{Verified}: \code{unittest discover} run \\
Empty-board opening is corner 0, not the centre & \textbf{Verified}: function called directly \\
Full empty-board search takes about 0.055\,s & \textbf{Verified}: timed on this machine \\
Test 3 duplicates test 1 & \textbf{Verified}: literals compared \\
\code{mark\_just\_played} is unused & \textbf{Verified}: body read \\
\code{after} handle discarded, no \code{after\_cancel} & \textbf{Verified}: grep \\
The resulting AI double-move & \textbf{Inferred}: traced through the source; no display available to drive the GUI \\
The 10 in \code{\_score\_terminal} must exceed max depth & \textbf{Inferred}: my reading of the arithmetic; unstated in the code \\
Console-version origin and its X-only win bug & \textbf{From the README}; the original script is not in this repository \\
Screenshot predates the AI mode & \textbf{Verified}: image viewed, no radio buttons present \\
\bottomrule
\end{tabular}
\caption{Where each non-obvious claim comes from.}
\end{table}

\end{document}
